\section{Introduction}

Recent advances in artificial intelligence and foundation models are driving a renewed interest in \emph{agentic systems} that can perceive, reason, and act in the world to complete tasks on our behalf \cite{russellnorvig1996, Wooldridge1995}. These systems promise to enhance our productivity by relieving us from mundane and laborious tasks, and revolutionize our lives by augmenting our knowledge and capabilities \cite{Jennings1998Applications,Wang2023survey, cheng2024exploring}. 
By leveraging the powerful reasoning and generative capabilities of large language models (LLMs), agentic systems are already making strides in fields like software engineering  \cite{yang2024swe,wang2024opendevinopenplatformai}, data analysis \cite{cao2024spider2vfarmultimodalagents}, scientific research \cite{lu2024ai,d2024marg} and web navigation \cite{zhou2024webarenarealisticwebenvironment,zhang2024webpilotversatileautonomousmultiagent}. 


Realizing the vision of agentic systems to transform our lives requires these systems to not only achieve high performance in specific domains, but also to generalize to the diverse range of tasks people may encounter throughout their day-to-day work and personal lives. In this paper, we take steps towards creating such a \emph{generalist agentic system} by introducing \emph{\team}.\footnote{The name \team{} is a combination of the words \textit{\textbf{m}ulti} and \textit{\textbf{agentic}}.} \team uses a team ofagents, each specializing in generally-useful skills, such as: operating a web browser, handling files, and executing code. The team is directed by an Orchestrator agent which guides progress towards a high-level goal by iteratively planning, maintaining working memory of progress, assigning tasks to other agents, and retrying upon encountering errors. The Orchestrator uses two \emph{structured ledgers} to achieve this and also to decide which agent should take the next action. 
Together, \team{}'s agents achieve strong performance on multiple challenging agentic benchmarks. 
Figure \ref{fig:magentic-example} shows an example of \team solving one such benchmark task that requires multiple steps and diverse tools. 

Key to \team's performance is its modular and flexible multi-agent approach \cite{MASaSurvey2000,Messing2002AnIT,Tambe1998AgentTeams,guo2024large,talebirad2023multiagentcollaborationharnessingpower},  implemented via the AutoGen\footnote{\url{https://github.com/microsoft/autogen}} framework \cite{wu2023autogen}. The multi-agent paradigm offers numerous advantages over monolithic single-agent approaches \cite{MASaSurvey2000,Tambe1998AgentTeams,cheng2024exploring,xi2023risepotentiallargelanguage}, which we believe makes it poised to become the leading paradigm in agentic development. For example, encapsulating distinct skills in separate agents simplifies development and facilitates reusability, akin to object-oriented programming. \team's specific design further supports easy adaptation and extensibility by enabling agents to be added or removed without altering other agents, or the overall workflow, unlike single-agent systems that often struggle with constrained and inflexible workflows.

To rigorously evaluate \team's performance, we introduce \emph{\agbench{}}, an extensible standalone tool for running agentic benchmarks. \agbench{}'s design enables repetition, isolation, and strong controls over initial conditions, so as to accommodate the variance of stochastic LLM calls, and to isolate the side-effects of agents taking actions.  Using \agbench{}, we evaluated \team on three agentic benchmarks. We observed task-completion rates of 38\% on GAIA \cite{mialon2023gaiabenchmarkgeneralai} and 32.8\% on WebArena \cite{zhou2024webarenarealisticwebenvironment}; and attained an accuracy of 27.7\% on AssistantBench \cite{yoran2024assistantbenchwebagentssolve}. These results place  \team in a strong position, where it is statistically competitive with other state-of-the-art (SOTA) systems, including those that are specialized for a given benchmark.  
Follow-up ablation experiments and in-depth error analyses reveal the additive value of each agent to \team's performance, and highlight opportunities for further improvement. 

In summary, we contribute: 
\begin{enumerate}
    \item \emph{\team}, a generalist multi-agent team with an open-source implementation. The team consists of five agents: a Coder, Computer Terminal, File Surfer, Web Surfer, and Orchestrator. Different agents can operate relevant tools such as stateful Web and file browsers, as well as command line and Python code executors. The Orchestrator  performs several functions to guide progress towards accomplishing a high-level goal: it formulates a plan, maintains structured working memory of progress, directs tasks to other agents, restarts and resets upon stalling, and determines task completion. 
    \item \emph{\agbench{}}, a standalone tool for evaluating systems on agentic benchmarks, also made available open-source.\footnote{\url{https://aka.ms/agbench}} \agbench{} handles  configuring, running, and reporting performance of agentic solutions while ensuring that all experiments start with well-known initial conditions, and that agents cannot interfere with one another across runs. 
    \item Experimental results and analyses of \team's performance on the GAIA, WebArena, and AssistantBench benchmarks, demonstrating strong task completion rates which are statistically competitive with other SOTA systems. We also examine the contribution of individual agents and capabilities, and provide an error analysis to identify the strengths and weaknesses of our multi-agent approach, along with opportunities for improvement.
\end{enumerate}